\documentclass[11pt]{article}
\usepackage[margin=2.4cm]{geometry}
\usepackage{amsmath,amssymb,bm,slashed}
\usepackage[colorlinks=true,linkcolor=blue,citecolor=blue,urlcolor=blue]{hyperref}
\title{QED radiative corrections for exclusive meson electroproduction
$ep\to e'\,B\,h$ ($h=\rho^0,\phi,J/\psi$ or $\pi^0,\pi^+,\eta$) with the
full set of structure functions
$\sigma_T,\ \sigma_L,\ \sigma_{TT},\ \sigma_{LT},\ \sigma_{LT'}$\\[2mm]
\large Formalism and description of the \texttt{vmrad} program}
\author{}
\date{\today}

\newcommand{\Li}{\mathrm{Li}_2}
\begin{document}
\maketitle

\section{Motivation}

The program DIFFRAD \cite{diffrad} computes QED radiative corrections (RC)
to exclusive vector--meson electroproduction
\begin{equation}
e(k_1)+p(p)\ \longrightarrow\ e'(k_2)+p'(p')+V(p_v),
\qquad V=\rho^0,\ \phi,\ J/\psi ,
\end{equation}
using only the two structure functions $\sigma_T$ and $\sigma_L$ of the
$\phi$-integrated cross section. The RC, however, depends on the structure
functions (SFs) used as input: the radiative ``tail'' entering a given
observed bin is built from the cross section at \emph{shifted} (true)
kinematics, and the azimuthal modulations generated by
$\sigma_{TT},\sigma_{LT},\sigma_{LT'}$ both receive and induce corrections
that a two-SF calculation cannot describe. Here we write down the RC for the
complete set of five SFs, including beam polarization, in the exact
covariant approach of Bardin and Shumeiko \cite{bs} as implemented for
exclusive pion electroproduction in EXCLURAD \cite{exclurad}, adapted to
vector--meson production. As in DIFFRAD/EXCLURAD, radiation from the lepton
line, the lepton vertex correction and vacuum polarization are taken into
account; radiation from the hadronic line and two--photon exchange are
neglected.

\section{Kinematics and conventions}

Metric $(+,-,-,-)$; $m$ is the lepton mass, $M$ the proton mass, $m_V$ the
meson mass. Observed (reconstructed from $e'$ and $V$) variables:
\begin{equation}
q=k_1-k_2,\quad Q^2=-q^2,\quad W^2=(p+q)^2,\quad t=(q-p_v)^2,\quad \phi,
\end{equation}
with $\phi$ the azimuth of the hadron plane about $\bm q$ measured from the
lepton plane. All four--vectors are constructed in the CM frame of $(p+q)$:
$\bm q$ along $+z$, leptons in the $x$--$z$ plane with $k_{1x}>0$,
$y=z\times x$ (Trento--like). Invariants:
\begin{equation}
S=2k_1\!\cdot\!p,\quad X=2k_2\!\cdot\!p = S-S_x,\quad S_x=2p\!\cdot\!q=W^2+Q^2-M^2,
\end{equation}
\begin{equation}
\lambda_S=S^2-4m^2M^2,\qquad
\lambda_q=S_x^2+4M^2Q^2,\qquad y=S_x/S .
\end{equation}
The virtual--photon polarization parameter (massless--lepton form, used only
for quoting the canonical decomposition; the code keeps $m$ exactly)
\begin{equation}
\varepsilon=\frac{1-y-\tfrac14\gamma^2y^2}{1-y+\tfrac12 y^2+\tfrac14\gamma^2y^2},
\qquad \gamma^2=\frac{4M^2x_B^2}{Q^2},\quad x_B=\frac{Q^2}{S_x}.
\end{equation}

For the radiative process $ep\to e'p'V\gamma(k)$ the reconstructed missing
mass moves away from the proton mass. We use the \emph{inelasticity}
\begin{equation}
v \equiv \Lambda^2-M^2,\qquad \Lambda\equiv p+q-p_v = p'+k ,
\label{eq:v}
\end{equation}
so $v=M_X^2-M^2$: an experimental missing--mass cut $M_X^2<M^2_{X,\rm cut}$
is exactly a cut $v<v_{\rm cut}$. At the Born level $v=0$. At fixed
$(Q^2,W^2,t)$ the meson CM energy is
$E_v=(W^2+m_V^2-M^2-v)/2W$ and
\begin{equation}
\cos\theta_v=\frac{t+Q^2-m_V^2+2q_0E_v}{2|\bm q||\bm p_v|},
\end{equation}
which yields the kinematic limit $v^{\rm kin}_{\max}(Q^2,W^2,t)$ from
$|\cos\theta_v|\le1$ (bounded globally by $(W-m_V)^2-M^2$). The integration
limit is $v_m=\min(v^{\rm kin}_{\max},\,v_{\rm cut})$.

\section{Born cross section and structure functions}

The Born matrix element is $\mathcal M_0=(e^2/Q^2)\,j^0_\mu H^\mu$ with
$j^0_\mu=\bar u(k_2)\gamma_\mu u(k_1)$. With the exact three--body phase
space one finds
\begin{equation}
\boxed{\
\frac{d\sigma_0}{dQ^2\,dW^2\,dt\,d\phi}
=\frac{\alpha^2}{32\pi^2\,\lambda_S\,\lambda_q^{1/2}\,Q^4}\;
L^0_{\mu\nu}H^{\mu\nu} \ }
\label{eq:born}
\end{equation}
where $L^0_{\mu\nu}(h)=\sum_{\rm spins}j^0_\mu j^{0*}_\nu$ for beam helicity
$h=\pm1$, and $H^{\mu\nu}=\overline{\sum}H^\mu H^{\nu*}$ is the hadronic
tensor of $\gamma^*p\to Vp'$.

\paragraph{Structure function definitions.}
Let $\ell^\mu$ (longitudinal) and $X^\mu,Y^\mu$ (transverse, $X$ in the
hadron plane, $Y$ normal to it) be the photon polarization basis in the
$\gamma^*p$ CM. The SFs are Hand--convention virtual--photoproduction cross
sections,
\begin{equation}
\frac{d\sigma_{\lambda\lambda'}}{dt}
=\frac{e^2\ \epsilon_\lambda^{\mu*}H_{\mu\nu}\epsilon_{\lambda'}^{\nu}}
{16\pi\,(W^2-M^2)^2},
\end{equation}
in terms of which the hadronic tensor reads
\begin{equation}
\boxed{\ H^{\mu\nu}=\frac{\kappa}{2}\Big[
\sigma_T\big(X^\mu X^\nu{+}Y^\mu Y^\nu\big)
+\sigma_L\,\ell^\mu \ell^\nu
+\sigma_{TT}\big(X^\mu X^\nu{-}Y^\mu Y^\nu\big)
-\sigma_{LT}\big(\ell^\mu X^\nu{+}X^\mu \ell^\nu\big)
-i\,\sigma_{LT'}\big(\ell^\mu X^\nu{-}X^\mu \ell^\nu\big)
\Big]\ }
\label{eq:htensor}
\end{equation}
with $\kappa=8(W^2-M^2)^2/\alpha$
(each $\sigma_i=d\sigma_i/dt$). Inserting \eqref{eq:htensor} into
\eqref{eq:born} reproduces \emph{identically} (this is verified to machine
precision in the code, for both beam helicities --- the ``Born closure''
test) the canonical decomposition
\begin{equation}
\frac{d\sigma_0}{dQ^2 dW^2 dt\,d\phi}
=\frac{\Gamma_v}{2\pi}\Big[\sigma_T+\varepsilon\,\sigma_L
+\varepsilon\cos2\phi\,\sigma_{TT}
+\sqrt{2\varepsilon(1{+}\varepsilon)}\cos\phi\,\sigma_{LT}
+h\sqrt{2\varepsilon(1{-}\varepsilon)}\sin\phi\,\sigma_{LT'}\Big].
\label{eq:canon}
\end{equation}
If your SF fits use $\sqrt{\varepsilon(1{+}\varepsilon)}$ instead of
$\sqrt{2\varepsilon(1{+}\varepsilon)}$, or opposite signs for
$\sigma_{TT},\sigma_{LT}$, rescale before plugging them in.

\section{Radiative corrections: master formula}

With the Bardin--Shumeiko covariant separation of the infrared divergence
(no soft--photon cutoff parameter appears anywhere),
\begin{equation}
\boxed{\ \sigma_{\rm obs}
=\sigma_0\,e^{\delta_{\rm inf}}\,(1+\delta_{VR}+\delta_{\rm vac})
+\sigma_F\ }
\label{eq:master}
\end{equation}
(EXCLURAD Eq.~(55) \cite{exclurad}); the exponentiation of
$\delta_{\rm inf}$ accounts for multiple soft photons.

\subsection{Virtual and soft parts}

Let $l_m=\ln(Q^2/m^2)$ and
\begin{equation}
S'_0=2k_1\!\cdot\!p'_B,\qquad X'_0=2k_2\!\cdot\!p'_B,\qquad
p'_B=p+q-p_v\big|_{v=0}\ \ (\text{Born recoil proton}).
\end{equation}
Then
\begin{align}
\delta_{\rm inf}&=\frac{\alpha}{\pi}\,(l_m-1)\,
\ln\frac{v_m^2}{S'_0X'_0},
\label{eq:dinf}\\[1mm]
\delta_{VR}&=\frac{\alpha}{\pi}\left(
\frac32 l_m-2-\frac12\ln^2\frac{X'_0}{S'_0}
+\Li\!\Big[1-\frac{Q^2M^2}{S'_0X'_0}\Big]-\frac{\pi^2}{6}\right).
\label{eq:dvr}
\end{align}
Note that $S'_0,X'_0$ depend on the hadron kinematics (through $p_v$), so
$\delta_{VR},\delta_{\rm inf}$ vary with $t$ and $\phi$ as well as with the
missing--mass cut through $v_m$.

Vacuum polarization: $\delta_{\rm vac}=\delta^{\rm lept}_{\rm vac}
+\delta^{\rm had}_{\rm vac}$ with the exact one--loop leptonic part
\begin{equation}
\delta^{\rm lept}_{\rm vac}
=2\sum_{\ell=e,\mu,\tau}\bar\Pi_\ell(Q^2),\qquad
\bar\Pi_\ell=\frac{2\alpha}{\pi}\int_0^1\!dx\,x(1-x)
\ln\!\Big[1+\frac{Q^2x(1-x)}{m_\ell^2}\Big],
\end{equation}
and the standard hadronic parameterization
$\delta^{\rm had}_{\rm vac}=-2\big(a+b\ln(1+cQ^2)\big)$ with the usual
piecewise coefficients (as in the POLRAD/HAPRAD family).

\subsection{Hard--photon (finite) part $\sigma_F$}

For the radiative process the true virtual photon is $\tilde q=q-k$ and the
hadronic vertex is evaluated at the \emph{shifted} kinematics
\begin{equation}
\tilde Q^2=Q^2+2k\!\cdot\!q,\qquad
\tilde W^2=W^2-2k\!\cdot\!(p+q),\qquad
\tilde t=t-2k\!\cdot\!(q-p_v).
\label{eq:shifted}
\end{equation}
Eliminating the unobserved proton with
$\delta^{(4)}$ and its mass shell, the photon phase space reduces, at fixed
observed $(Q^2,W^2,t,\phi)$, to an integral over $v$ and the photon
direction $\bm n$: with
\begin{equation}
z\equiv \frac{k\cdot\Lambda}{\omega}=\Lambda_0-\bm\Lambda\!\cdot\!\bm n,
\qquad \omega=\frac{v}{2z},
\end{equation}
the exact radiative cross section is
\begin{equation}
\frac{d\sigma_R}{dQ^2dW^2dt\,d\phi}
=\frac{\alpha^3}{512\,\pi^4\,\lambda_S\,\lambda_q^{1/2}}
\int_0^{v_m}\!dv\int d\Omega_k\;
\frac{v}{z^2}\;
\frac{L^R_{\mu\nu}\,\tilde H^{\mu\nu}}{\tilde Q^4},
\label{eq:sigR}
\end{equation}
where $\tilde H^{\mu\nu}$ is the tensor \eqref{eq:htensor} built from
$\tilde q$, $p$, $p_v(v)$ with all five SFs evaluated at
$(\tilde W^2,\tilde Q^2,\tilde t)$, and $L^R_{\mu\nu}(h)$ is the exact
radiative leptonic tensor
\begin{equation}
L^R_{\mu\nu}=\sum_{\rm spins,\ pol}j^R_\mu j^{R*}_\nu,\qquad
j^R_\mu=\bar u(k_2)\Big[\gamma_\mu\frac{\slashed k_1-\slashed k+m}{-\varkappa_1}
\slashed\epsilon^*
+\slashed\epsilon^*\frac{\slashed k_2+\slashed k+m}{\varkappa_2}\gamma_\mu
\Big]u(k_1,h),
\end{equation}
$\varkappa_{1,2}=2k\!\cdot\!k_{1,2}$. In the soft limit
$L^R\to 4F_{\rm IR}L^0$ with the eikonal factor
\begin{equation}
F_{\rm IR}=\frac{Q^2+2m^2}{\varkappa_1\varkappa_2}
-\frac{m^2}{\varkappa_1^2}-\frac{m^2}{\varkappa_2^2}.
\end{equation}
The IR divergence of \eqref{eq:sigR} at $v\to0$ is removed covariantly by
subtracting the eikonal times the \emph{unshifted} Born integrand --- whose
analytic integral is precisely what produced
\eqref{eq:dinf}--\eqref{eq:dvr} --- leaving the finite
\begin{equation}
\boxed{\ \sigma_F=\frac{\alpha^3}{512\pi^4\lambda_S\lambda_q^{1/2}}
\int_0^{v_m}\!\!dv\!\int\! d\Omega_k\,\frac{v}{z^2}
\left[\frac{L^R_{\mu\nu}\tilde H^{\mu\nu}}{\tilde Q^4}
-\,4F_{\rm IR}\,\frac{L^0_{\mu\nu}H^{0\,\mu\nu}}{Q^4}\right] \ }
\label{eq:sigF}
\end{equation}
(the analogue of EXCLURAD Eq.~(38)). Every azimuthal interference structure
of all five SFs at shifted kinematics is generated \emph{exactly} by the
tensor contraction $L^R\!:\!\tilde H$ --- no leading--log or peaking
approximation is made.

\section{Numerical implementation (\texttt{vmrad})}

\begin{itemize}
\item $L^0,L^R$ are evaluated \emph{numerically} from helicity spinors
(exact in $m$, exact beam polarization), instead of transcribing the very
long analytic contractions of \cite{exclurad}. Correctness is enforced by
the validation suite: Born closure against \eqref{eq:canon} ($\sim3\times
10^{-8}$, limited only by lepton--mass terms), current conservation
$\tilde q_\mu L^{R\,\mu\nu}=0$ ($\sim10^{-9}$), and soft factorization
$L^R\to4F_{\rm IR}L^0$.
\item The polarization basis $(X,Y,\ell)$ at shifted kinematics is built by
boosting to the $(p+\tilde q)$ CM; the SFs enter as scalars
$\sigma_i(\tilde W^2,\tilde Q^2,\tilde t)$ --- no shifted azimuth is ever
needed.
\item The photon angular integral uses a two--channel decomposition with
weights $w_c\propto\varkappa_c^{-2}$ centered on $\bm k_1$ and $\bm k_2$ and
the substitution $x=\ln(a_c-\cos\theta_c)$, which flattens the collinear
peaks (width $m^2/E^2$); the collinear core gets a dedicated
Gauss--Legendre panel. The $v$ integral is Gauss--Legendre in $\ln v$.
\item IR bookkeeping is verified non--circularly: the angular integral of
the subtraction term in \eqref{eq:sigF} reproduces the analytic slope
$d\,\delta_{\rm inf}/d\ln v^2=(\alpha/\pi)(l_m-1)$ to $\sim2\times10^{-6}$.
\item Typical grids ($n_v=30$, $16{+}24$ polar, $16$ azimuthal nodes per
channel) give $\sigma_F$ to $\sim10^{-4}$ relative.
\end{itemize}

\paragraph{Inputs and outputs.}
The user supplies $\sigma_i(W^2,Q^2,t)$ in nb/GeV$^2$ (class
\texttt{StructureFunctions}; demo $\rho^0/\phi/J\psi$ models included, to be
replaced for physics). \texttt{rc\_factor(...)} returns $\sigma_{\rm Born}$,
$\sigma_{\rm obs}$, $\eta=\sigma_{\rm obs}/\sigma_{\rm Born}$ and all pieces
of \eqref{eq:master}, at fixed $(E_{\rm beam},Q^2,W^2,t,\phi,h)$ and
missing--mass cut. A \texttt{TwoSF} wrapper zeroes
$\sigma_{TT},\sigma_{LT},\sigma_{LT'}$ for direct comparison with the
DIFFRAD--like two--SF treatment. Because the SFs are sampled at
\eqref{eq:shifted}, the model must be defined down to threshold in
$\tilde W^2$ and for $\tilde Q^2$ both above and below the observed $Q^2$.

\section{Pseudoscalar mesons}

Nothing in Secs.~2--5 refers to the spin of the produced meson: the
five--SF decomposition \eqref{eq:htensor}--\eqref{eq:canon} is precisely the
one introduced for exclusive \emph{pion} electroproduction (it is
EXCLURAD's original case), so the identical formulas and code apply to
$ep\to e'p\pi^0$, $e'n\pi^+$, $e'p\eta$, with two mass substitutions:
\begin{itemize}
\item the detected--meson mass $m_V\to m_h$ ($m_{\pi^0},m_{\pi^+},m_\eta$)
in the hadron kinematics and in $t_{\min}$;
\item the \emph{undetected recoil--baryon} mass $m_U$ (proton for
$\pi^0 p,\ \eta p$ and the vector channels; \emph{neutron} for $\pi^+n$)
everywhere the missing system appears: in the inelasticity
$v=\Lambda^2-m_U^2=M_X^2-m_U^2$ and its limits
(globally $v\le (W-m_h)^2-m_U^2$), and in the $\Li$ argument of
$\delta_{VR}$, Eq.~\eqref{eq:dvr}, which becomes
$\Li[1-Q^2m_U^2/(S'_0X'_0)]$ (it originates from the soft--photon integral
in the rest frame of $\Lambda$, whose mass at $v\to0$ is $m_U$).
\end{itemize}
In \texttt{vmrad} the SF model simply carries \texttt{mV} and \texttt{mU}
attributes (demo models \texttt{demo\_pi0}, \texttt{demo\_pip},
\texttt{demo\_eta} are provided); the Born closure and IR validation tests
pass identically for the pseudoscalar case, including $m_U\ne M$
($\pi^+n$). For $\pi^+$ production the emission from the charged meson
line (and from the hadronic line in general) remains neglected, exactly as
in EXCLURAD.

\paragraph{Limitations.} Lepton--line radiation only (hadronic radiation
and two--photon exchange neglected, as in DIFFRAD/EXCLURAD);
$\delta^{\rm had}_{\rm vac}$ from a fit; exclusivity is assumed perfect
below the missing--mass cut (no detector response). The correction is
defined for the unshifted--bin observable
$d\sigma/dQ^2dW^2dt\,d\phi$ with $t$ reconstructed as $(q-p_v)^2$; if your
analysis reconstructs $t$ differently (e.g.\ from the recoil proton), the
shifted variables \eqref{eq:shifted} must be adapted.

\begin{thebibliography}{9}
\bibitem{diffrad} I.~Akushevich,
``Radiative effects in processes of diffractive vector meson
electroproduction,'' Eur.\ Phys.\ J.\ C \textbf{8} (1999) 457
[hep-ph/9808309]; see also I.~Akushevich and P.~Kuzhir,
``QED radiative correction to spin--density matrix elements in exclusive
vector meson production,'' Phys.\ Lett.\ B \textbf{474} (2000) 411
[hep-ph/9911346].
\bibitem{bs} D.~Yu.~Bardin and N.~M.~Shumeiko,
Nucl.\ Phys.\ B \textbf{127} (1977) 242.
\bibitem{exclurad} A.~Afanasev, I.~Akushevich, V.~Burkert and K.~Joo,
``QED radiative corrections in processes of exclusive pion
electroproduction,'' Phys.\ Rev.\ D \textbf{66} (2002) 074004
[hep-ph/0208183]; code EXCLURAD, \url{https://www.jlab.org/RC/}.
\end{thebibliography}

\end{document}
